\documentclass[11pt,a4paper]{article}

% --- Math, symbols, theorem environments --------------------------------
\usepackage{amsmath,amssymb,amsthm,mathtools}

% --- Fonts --------------------------------------------------------------
% TeX Gyre Pagella (Palatino clone) for body text;
% Monaspace Xenon for monospace, scaled to harmonize with Pagella;
% Pagella math companion via newpxmath (Palatino-flavoured math).
\usepackage{fontspec}
\setmainfont{TeX Gyre Pagella}
\setmonofont{Monaspace Xenon}[
  Scale=0.875,
  Ligatures=NoCommon,
  BoldFont={Monaspace Xenon SemiBold},
  ItalicFont={Monaspace Xenon Italic},
  BoldItalicFont={Monaspace Xenon SemiBold Italic},
]
% Sans-serif Monaspace variant used for math operator names
% (\eval, \depth, \embed, \largest).  Distinct from both the Pagella
% prose and the Xenon listings.
\newfontfamily{\monoargon}{Monaspace Argon}[
  Scale=0.875,
  Ligatures=NoCommon,
]
\usepackage{newpxmath}                  % Palatino-flavoured math symbols
\linespread{1.05}                       % Pagella benefits from a touch more leading
% Define [[ ]] as composite glyphs with a small *positive* gap between
% the two halves, so they are clearly readable as double brackets at
% any zoom level (unlike the negative-thin-space version, which can
% read as bold-single-brackets at small sizes).  We avoid [stmaryrd]
% because its [stmary10] font is not bundled with [lualatex] by
% default.
\providecommand{\llbracket}{\mathopen{[\mkern-3.25mu[}}
\providecommand{\rrbracket}{\mathclose{]\mkern-3.25mu]}}

% --- Layout, typography -------------------------------------------------
\usepackage[a4paper,margin=1in]{geometry}
\usepackage{microtype}
\usepackage{parskip}
\usepackage{enumitem}
\usepackage{booktabs}
\usepackage{xcolor}

% --- Code listings ------------------------------------------------------
\usepackage{listings}

% Keyword color #0E2D62 (HSL 217/75/22) and numeric-literal color
% #404D1A (HSL 76/50/20) sit at similar lightness so they read as
% paired highlights rather than competing.  The olive cast on
% literals is mainly there to make digit "1" unmistakable from
% letter "l" in the Monaspace Xenon body.
\definecolor{coqkw}{HTML}{0E2D62}     % vernacular, tactics, reserved
\definecolor{coqnum}{HTML}{404D1A}    % numeric literals (0..9)
\definecolor{coqcomment}{RGB}{105,120,140}
\definecolor{coqstr}{RGB}{160,40,40}
\definecolor{coqid}{RGB}{0,0,0}
\definecolor{coqbg}{RGB}{248,248,250}

\lstdefinelanguage{Coq}{
  morekeywords={Definition,Fixpoint,Inductive,Lemma,Theorem,Corollary,
                Proof,Qed,Defined,forall,exists,fun,match,with,end,let,in,
                if,then,else,return,as,struct,Module,End,Import,Require,
                From,Type,Prop,Set,Section,Variable,Hypothesis,Notation,
                Local,Global,Print,Eval,Compute,Time,where,Parameter,Axiom,
                Ltac,intro,intros,apply,eapply,exact,rewrite,erewrite,
                simpl,cbn,cbv,beta,delta,iota,zeta,eta,
                destruct,induction,reflexivity,symmetry,transitivity,
                assumption,auto,eauto,lia,split,exists,eexists,
                subst,assert,by,
                pose,proof,unfold,fold,change,replace,specialize,
                generalize,revert,clear,case,solve,contradiction,
                inversion,discriminate,congruence,repeat,first,try,now,
                progress,context,goal,at,all,
                projT1,projT2,existT,ex_intro,List,In,Some,None,nth_error,
                Nat,iter,recursion,nil,cons,nat,list,max,bool,option,True,False,
                Bz,Bsucc,Blim,FGH,BigGrow,epsilon_0,Record,
                tVar,tLam,tApp,tO,tS,tNatRec,tGrow,tpNat,tpArr,
                term,interp_type,pack,RelVal,RelPack,REntry,
                contender_5,contender_6},
  morecomment=[s]{(*}{*)},
  morestring=[b]",
  sensitive=true,
  alsoletter={'},
  basicstyle=\ttfamily\fontsize{10.5pt}{12.6pt}\selectfont,
  keywordstyle=\color{coqkw},
  commentstyle=\color{coqcomment}\itshape,
  stringstyle=\color{coqstr},
  identifierstyle=\color{coqid},
  backgroundcolor=\color{coqbg},
  frame=single,
  framerule=0.3pt,
  rulecolor=\color{coqcomment!40},
  xleftmargin=1em,
  xrightmargin=0.5em,
  framexleftmargin=0.5em,
  showstringspaces=false,
  breaklines=true,
  keepspaces=true,
  columns=fullflexible,
  literate=
    {<=}{{$\leq$}}2
    {>=}{{$\geq$}}2
    {<>}{{$\neq$}}2
    {<-}{{$\leftarrow\,$}}2
    {->}{{$\rightarrow$}}2
    {=>}{{$\Rightarrow$}}2
    {|-}{{$\vdash$}}2
    {/\\}{{$\land$}}2
    {\\/}{{$\lor$}}2
    {<->}{{$\leftrightarrow$}}3
    {Contender.}{{\textcolor{gray!55}{Contender.}}}{10}
    {LGrow.}{{\textcolor{gray!55}{LGrow.}}}{6}
    {forall}{{$\forall$}}1
    {exists}{{$\exists$}}1
    {0}{{\textcolor{coqnum}{0}}}1
    {1}{{\textcolor{coqnum}{1}}}1
    {2}{{\textcolor{coqnum}{2}}}1
    {3}{{\textcolor{coqnum}{3}}}1
    {4}{{\textcolor{coqnum}{4}}}1
    {5}{{\textcolor{coqnum}{5}}}1
    {6}{{\textcolor{coqnum}{6}}}1
    {7}{{\textcolor{coqnum}{7}}}1
    {8}{{\textcolor{coqnum}{8}}}1
    {9}{{\textcolor{coqnum}{9}}}1
    % Passthrough rule(s) below: listings tries last-declared literate
    % patterns first.  Listing identifiers that share a prefix with
    % "forall" or "exists" (which are substituted above) need to appear
    % here verbatim so the longer match wins and the prefix-substitution
    % is never applied.  Add new entries as new identifiers appear.
    {exists_maximizer_42}{{exists\textunderscore{}maximizer\textunderscore{}\textcolor{coqnum}{4}\textcolor{coqnum}{2}}}{19}
}

\lstset{language=Coq}

% --- Theorem environments -----------------------------------------------
\theoremstyle{plain}
\newtheorem{theorem}{Theorem}[section]
\newtheorem{lemma}[theorem]{Lemma}
\newtheorem{proposition}[theorem]{Proposition}
\newtheorem{corollary}[theorem]{Corollary}

\theoremstyle{definition}
\newtheorem{definition}[theorem]{Definition}
\newtheorem{example}[theorem]{Example}

\theoremstyle{remark}
\newtheorem{remark}[theorem]{Remark}

% --- Math shorthands ----------------------------------------------------
\newcommand{\Nat}{\mathbb{N}}
\newcommand{\eps}{\varepsilon}
\newcommand{\opname}[1]{\mathop{\text{\monoargon #1}}\nolimits}
\newcommand{\depth}{\opname{depth}}
\newcommand{\eval}{\opname{eval}}
\newcommand{\embed}{\opname{embed}}
\newcommand{\largest}{\opname{largest}}

% --- Hyperref last ------------------------------------------------------
\usepackage[colorlinks=true,
            linkcolor=blue!50!black,
            citecolor=blue!50!black,
            urlcolor=blue!50!black]{hyperref}

% --- Title & metadata ---------------------------------------------------
\title{From Ackermann to $\varepsilon_0$:\\
       a Tutorial Walk-Through of \texttt{contender\_6}}
\author{(name-the-biggest-number, contender~6)}
\date{}

\begin{document}
\maketitle
\thispagestyle{empty}

\begin{abstract}
\noindent
This article is a self-contained tutorial on the construction of
\texttt{contender\_6}, a new entry to the Coq-formalized
\emph{name-the-biggest-number} contest of
\texttt{github.com/codyroux/name-the-biggest-number}.  The previous
contender, \texttt{contender\_5}, is the largest natural number that any
closed term of System~T (simply-typed $\lambda$-calculus with
\texttt{NatRec}) of abstract-syntax-tree depth at most $42$ evaluates to.
We extend that language with a single new primitive, \texttt{tGrow},
whose semantics is the function $f_{\varepsilon_0}$ from the fast-growing
hierarchy, and define \texttt{contender\_6} as the analogous depth-bounded
maximum at depth $44$ in the extended language.  The proof of the strict
inequality $\texttt{contender\_5} < \texttt{contender\_6}$ embeds the
previous contender's (existentially obtained) maximizer term into the
new language and applies \texttt{tGrow} on top of it, yielding the
explicit lower bound
$\texttt{contender\_6} \geq f_{\varepsilon_0}(\texttt{contender\_5}+1)$.

The tutorial presents, in turn: the contest setting and its
\emph{don't-be-lazy} rule; the machinery behind \texttt{contender\_5}
that we inherit; Brouwer ordinal notations, which turn ordinal
recursion into ordinary structural recursion; the fast-growing
hierarchy; the new object language $\mathcal{L}_{\textrm{Grow}}$; and
the three-ingredient proof, whose only subtle part is a Tait-style
logical relation showing that the embedding preserves evaluation
\emph{without} assuming function extensionality.  Every step is
accompanied by the corresponding Rocq/Coq code, taken verbatim from
\texttt{Contender.v}; the whole development is axiom-free.
\end{abstract}

\tableofcontents
\newpage

% =======================================================================
\section{The contest in a nutshell}
% =======================================================================

\subsection{Big numbers, formally}

In Scott Aaronson's 1999 essay
\emph{Who Can Name the Bigger Number?}~\cite{aaronson} the rules of
the eponymous game are a single line on a blackboard: each contestant
writes a natural number, and the larger one wins.  The interesting twist
is that ``larger'' is not measured in digits or magnitude of the numeral,
but in the \emph{description} of the number.  Tower notation beats
exponentials, primitive recursion beats towers, the Ackermann function
beats primitive recursion, and so on, climbing through the ordinals of
proof theory.

The Coq repository \texttt{codyroux/name-the-biggest-number} formalizes
this game in a small but precise way:

\begin{itemize}[leftmargin=2em]
  \item Numbers must be \emph{constructive}: no busy-beaver shenanigans,
        no oracles, no inconsistent axioms.  Each contender is a closed
        term of type \texttt{nat}.
  \item Each contender comes with a \emph{formal Coq proof} that it
        strictly exceeds the previous one.
  \item Type-checking the definition takes at most $15$ seconds and
        the proof at most $60$ seconds on a reasonable machine.
  \item No axioms.
  \item The crucial \emph{don't-be-lazy} rule: the new contender
        \texttt{contender\_N} must \emph{not} be defined by invoking the
        previous \texttt{contender\_N-1}.  The spirit of the rule is
        that a one-line victory of the form ``$\texttt{contender\_5} +
        1$'' should be ineligible.  Section~\ref{sec:why-c5-plus-1}
        discusses what this means for us in practice.
\end{itemize}

\subsection{The chain so far}

\begin{center}
\begin{tabular}{cl@{\ }l}
\toprule
\# & Definition & Order of magnitude \\
\midrule
1 & $42$ & a two-digit literal \\
2 & a tower of $42$\,\textasciicircum\,$\cdots$\,\textasciicircum\,$42$ & several thousand $42$s tall \\
3 & $\texttt{tower\_of\_power}\ 10000$ & a $42$-tower of height $10001$ \\
4 & $\texttt{ack}\ 5\ 42\ 10002$ & an Ackermannian giant \\
5 & largest closed System~T term, $\depth \leq 42$ & dwarfs every $f_\alpha$ for $\alpha < \varepsilon_0$ \\
\textbf{6} & \textbf{the present submission} & $\geq f_{\varepsilon_0}(\texttt{contender\_5}+1)$ \\
\bottomrule
\end{tabular}
\end{center}

The first four contenders climb the classical ladder -- a literal, a
tower of exponentials, a ten-thousand-story tower built by primitive
recursion, and an Ackermann value one hierarchy step above that.  Their
definitions and proofs live in the upstream repository
(\texttt{Graveyard.v} and \texttt{Contender.v}) and play no role in
what follows: our construction interacts \emph{only} with
\texttt{contender\_5}.  We therefore recap just that one, in exactly
the amount of detail the new proof needs.

% =======================================================================
\section{The previous contender: a depth-bounded maximum over System~T}
\label{sec:contender-5}
% =======================================================================

\texttt{contender\_5} was the qualitative leap in the chain.  Instead
of naming a specific large number, it names \emph{the largest number
that any closed System~T term of depth at most $42$ evaluates to}.
This section walks through the machinery that realizes this idea in
Coq, because \texttt{contender\_6} runs the very same recipe -- in a
bigger language.

\subsection{The object language}

The language is simply-typed $\lambda$-calculus over one base type,
extended with the natural numbers and their recursor -- G\"odel's
System~T, in modern clothing.  Types and their interpretations as Coq
types are:

\begin{lstlisting}
Inductive type :=
| tpNat : type
| tpArr : type -> type -> type.

Fixpoint interp_type (tp : type) : Type :=
  match tp with
  | tpNat       => nat
  | tpArr t1 t2 => interp_type t1 -> interp_type t2
  end.
\end{lstlisting}

Terms are de~Bruijn-indexed, with the slightly unusual convention
that index $0$ is the \emph{outermost} binder (de~Bruijn
\emph{levels} rather than indices):

\begin{lstlisting}
Inductive term :=
| tVar    (x : nat)
| tLam    (A B : type) (body : term)
| tApp    (t1 t2 : term)
| tO
| tS
| tNatRec (R : type).
\end{lstlisting}

Note that terms carry \emph{type annotations} but no typing
\emph{derivation}: \texttt{termsUpTo} below will enumerate plenty of
ill-typed terms such as $\texttt{tApp}\;\texttt{tO}\;\texttt{tO}$.
The interpreter handles them with a default-value strategy.  It maps a
term and an environment of values to a \emph{pack}: a dependent pair
of a type tag and a value of that type,

\begin{lstlisting}
Definition interp_term :
    forall (e : list {tp : type & interp_type tp})
           (t : term),
      {tp : type & interp_type tp}.
\end{lstlisting}

with the following behavior.  A variable is looked up in the
environment (\texttt{lookup} returns an \texttt{option}); a miss
produces the default pack $(\texttt{tpNat}, 0)$.  A $\lambda$ produces
a function pack at tag $\texttt{tpArr}\,A\,B$, coercing the
interpreted body to $B$ via a helper \texttt{cast}: a coercion between
any two type tags that is the identity when the tags agree and falls
back to a default \texttt{error} value (i.e.\ $0$ at \texttt{tpNat},
the constant-\texttt{error} function at arrow types) when they do not.
An application inspects the type tag of the function part: if it is an
arrow, the argument is cast to the domain and the function applied;
otherwise the result is the error pack.  Finally \texttt{tO},
\texttt{tS} and \texttt{tNatRec R} denote $0$, the successor, and the
standard library recursor \texttt{Nat.recursion} at return type
\texttt{R}.

The top-level evaluator extracts a natural number from a closed term,
returning the value at tag \texttt{tpNat} and $0$ at arrow tags:

\begin{lstlisting}
Definition eval (t : term) : nat :=
  let (tp, res) := interp_term nil t in
  match tp as t0 return interp_type t0 -> nat with
  | tpNat       => fun n => n
  | tpArr _ _   => fun _ => 0
  end res.
\end{lstlisting}

\subsection{Depth, and why there are only finitely many terms}

Term depth is abstract-syntax-tree depth, with two twists that make
depth-bounded enumeration possible: a variable index $x$ costs
$x+1$ (so only finitely many variables fit in any budget), and type
annotations count too (so only finitely many annotations fit):

\begin{lstlisting}
Definition nat_depth (n : nat) := S n.

Fixpoint type_depth (t : type) : nat :=
  match t with
  | tpNat       => 1
  | tpArr t1 t2 => S (max (type_depth t1) (type_depth t2))
  end.

Fixpoint term_depth (t : term) : nat :=
  match t with
  | tVar x        => S (nat_depth x)
  | tLam A B body => S (max (max (type_depth A)
                                 (type_depth B))
                            (term_depth body))
  | tApp t1 t2    => S (max (term_depth t1) (term_depth t2))
  | tO            => 1
  | tS            => 1
  | tNatRec R     => S (type_depth R)
  end.
\end{lstlisting}

``Finitely many'' is realized as concrete enumerator functions --
lists that provably contain everything of depth at most $n$:

\begin{lstlisting}
Fixpoint natsUpTo (n : nat) : list nat :=
  match n with
  | 0    => []
  | S m  => m :: natsUpTo m
  end.

Fixpoint typesUpTo (n : nat) : list type :=
  match n with
  | 0    => []
  | S m  =>
    let r := typesUpTo m in
    tpNat :: List.map (fun '(t1, t2) => tpArr t1 t2)
                      (list_prod r r)
  end.

Fixpoint termsUpTo (n : nat) : list term :=
  match n with
  | 0    => []
  | S m  =>
    List.map tVar (natsUpTo m) ++
    List.map (fun '(A, B, body) => tLam A B body)
             (list_prod (list_prod (typesUpTo m) (typesUpTo m))
                        (termsUpTo m)) ++
    List.map (fun '(t1, t2) => tApp t1 t2)
             (list_prod (termsUpTo m) (termsUpTo m)) ++
    [tO] ++ [tS] ++
    List.map tNatRec (typesUpTo m)
  end.
\end{lstlisting}

Each enumerator comes with a correctness lemma of the form
``membership in the list is \emph{equivalent} to fitting in the depth
budget'':

\begin{lstlisting}
Lemma natsUpTo_correct  : forall n m, nat_depth m  <= n <-> In m (natsUpTo n).
Lemma typesUpTo_correct : forall n t, type_depth t <= n <-> In t (typesUpTo n).
Lemma termsUpTo_correct : forall n t, term_depth t <= n <-> In t (termsUpTo n).
\end{lstlisting}

The proofs are routine but pleasant inductions; the term-level one has
one case per constructor, each navigating to the right \texttt{++}
segment.

\subsection{The argmax, and \texttt{contender\_5}}

The maximum-by-key combinator is an accumulator-style fold,
polymorphic in the element type:

\begin{lstlisting}
Fixpoint maxBy {T : Type}
               (f : T -> nat) (currentMax : nat)
               (currentBest : T) (l : list T) : T :=
  match l with
  | nil       => currentBest
  | cons h t  =>
    if currentMax <? f h
    then maxBy f (f h) h t
    else maxBy f currentMax currentBest t
  end.
\end{lstlisting}

Its two key properties, recorded upstream as Coq lemmas, are the
dichotomy ``the result is an element of the list or the initial
best'', and the lower-bound property ``every list element's key is at
most the result's key'':

\begin{lstlisting}
Lemma maxBy_In : forall f l currentMax currentBest,
    currentMax = f currentBest ->
    In (maxBy f currentMax currentBest l) l \/
    maxBy f currentMax currentBest l = currentBest.

Lemma lowerbound_maxBy : forall f x l currentMax currentBest,
    In x l ->
    f currentBest = currentMax ->
    f x <= f (maxBy f currentMax currentBest l).
\end{lstlisting}

And the number itself is the evaluation of the depth-$42$ argmax:

\begin{lstlisting}
Definition largest_of_depth (n : nat) : term :=
  maxBy eval 0 tO (termsUpTo n).

Definition largest_STLCNatRec_nat_of_depth (n : nat) : nat :=
  eval (largest_of_depth n).

Definition contender_5 : nat :=
  largest_STLCNatRec_nat_of_depth 42.
\end{lstlisting}

Of course nobody can \emph{run} this definition: \texttt{termsUpTo 42}
is a list of astronomical length, and evaluating its members includes
computing Ackermann-sized numbers and worse.  The definition is a
perfectly good closed term of type \texttt{nat} nonetheless, and all
reasoning about it goes through the lemmas above rather than through
computation.  (For the same reason, our development marks all of this
machinery \texttt{Opaque}, so that no Coq tactic accidentally asks the
kernel to normalize a depth-$42$ search.)

\subsection{What the new proof will reuse}

From everything above, the proof of
$\texttt{contender\_5} < \texttt{contender\_6}$ will need exactly
three inherited lemmas --

\begin{itemize}[leftmargin=2em]
  \item \texttt{termsUpTo\_correct} (in the direction ``member of
        $\texttt{termsUpTo}\,42$ $\Rightarrow$ depth $\leq 42$''),
  \item \texttt{maxBy\_In},
  \item \texttt{lowerbound\_maxBy},
\end{itemize}

-- plus the \texttt{interp\_term}/\texttt{eval} infrastructure itself
(\texttt{lookup}, \texttt{cast}, \texttt{error}), which the new
language shares.  The upstream proof of
$\texttt{contender\_4} < \texttt{contender\_5}$ -- an ingenious
reification of the Ackermann function as a System~T term -- is a
spectator: we never touch it.

% =======================================================================
\section{Why \texttt{contender\_5}$\,+\,$1 is not allowed}
\label{sec:why-c5-plus-1}
% =======================================================================

The naive way to defeat \texttt{contender\_5} is

\[
\texttt{contender\_6} \;\overset{?}{=}\; \texttt{contender\_5} + 1.
\]

This ``definition'' is exactly what the don't-be-lazy rule bans: the
\emph{definition} mentions \texttt{contender\_5}, even though the proof
would be one line.  And the ban is reasonably read as extending to any
definition whose unfolding contains \texttt{contender\_5}'s search
machinery (its \texttt{largest\_\dots\_of\_depth} wrappers, its
\texttt{interp\_term}, its \texttt{termsUpTo}) -- otherwise
\[
  \texttt{largest\_STLCNatRec\_nat\_of\_depth}\ 43
\]
would be a legal one-line submission, which is laziness in all but
name.

The rule is about \emph{definitions}, though, not \emph{proofs}.  A
proof of $\texttt{contender\_5} < \texttt{contender\_6}$ must of course
talk about \texttt{contender\_5}; the upstream proof of
$\texttt{contender\_4} < \texttt{contender\_5}$ likewise reasons about
\texttt{contender\_4} at length.  This gives us a clean design
principle:

\begin{quote}
\emph{The new contender's \textbf{definition} must not mention any of
the previous contender's machinery.  The \textbf{proof} may use
anything it likes -- including the fact that \texttt{contender\_5}'s
maximum is attained by some (unknowable) term.}
\end{quote}

% =======================================================================
\section{The plan: one fresh primitive}
\label{sec:plan}
% =======================================================================

System~T is famously tied to the ordinal $\varepsilon_0$: the
functions it can express are exactly those provably total in Peano
arithmetic, and every closed term of type \texttt{Nat} is bounded by
$f_\alpha(c)$ for some $\alpha < \varepsilon_0$ and some constant $c$.
So a function that \emph{diagonalizes over} that entire hierarchy --
the fast-growing-hierarchy function at $\varepsilon_0$ itself,

\[
   \texttt{BigGrow}(n) \;\coloneqq\; f_{\varepsilon_0}(n),
\]

-- is not the interpretation of any System~T term.  Adding it to the
language as a primitive therefore yields \emph{genuinely} new
expressive power, not a rearrangement of what depth $42$ could already
buy.  This suggests the whole construction in three lines:

\begin{enumerate}[leftmargin=2em,label=(\arabic*)]
  \item Define a fresh object language
        $\mathcal{L}_{\textrm{Grow}} = \text{System~T} \,+\, \texttt{tGrow}$,
        where \texttt{tGrow} is a new constant of type
        $\texttt{Nat}\to\texttt{Nat}$ interpreted as \texttt{BigGrow}.
  \item Define
        $\texttt{contender\_6} \coloneqq$ the largest $\eval(t)$ over
        closed $\mathcal{L}_{\textrm{Grow}}$ terms $t$ of depth at
        most $44$ -- the same argmax construction as before, two depth
        units richer.
  \item For the proof, obtain the \emph{existential} maximizer
        $t^\star$ of \texttt{contender\_5} (a closed System~T term
        with $\depth(t^\star)\leq 42$ and
        $\eval(t^\star)=\texttt{contender\_5}$, guaranteed to exist
        by \texttt{maxBy\_In}), embed it into
        $\mathcal{L}_{\textrm{Grow}}$, and offer the term
        \[
          \texttt{tGrow}\,\bigl(\texttt{tS}\,(\embed\,t^\star)\bigr)
        \]
        -- of depth at most $44$ -- as a witness to the new maximum.
\end{enumerate}

The witness evaluates to
$\texttt{BigGrow}(\texttt{contender\_5}+1)$, so the argmax at depth
$44$ is at least that, and $\texttt{BigGrow}(n+1) > n$ closes the
strict inequality.  Note how the two extra depth units are spent:
one on the \texttt{tS} (so that we exceed the old maximum even
\emph{before} \texttt{tGrow} acts), one on the \texttt{tGrow}.

The definition in step (2) mentions only the new language's own
evaluator and enumerator, plus the generic list combinator
\texttt{maxBy} -- no \texttt{contender\_5} machinery.  Everything
that touches \texttt{contender\_5} happens inside the proof, which
is exactly what the contest permits.

We must build \texttt{BigGrow} first, and it must be a bona fide
axiom-free total Coq function.  That is the business of the next two
sections, and it turns out to be surprisingly pleasant.

% =======================================================================
\section{Ordinals as data: Brouwer trees}
\label{sec:brouwer}
% =======================================================================

Readers familiar with the fast-growing hierarchy can skim this section
and the next; the Coq definitions are all four-liners.

\subsection{The type}

Brouwer ordinal notations represent countable ordinals as well-founded
trees:

\begin{lstlisting}
Inductive Brouwer : Type :=
| Bz    : Brouwer
| Bsucc : Brouwer -> Brouwer
| Blim  : (nat -> Brouwer) -> Brouwer.
\end{lstlisting}

\texttt{Bz} is the ordinal $0$; \texttt{Bsucc}\,$a$ is the successor
$a+1$; and \texttt{Blim}\,$f$ is the limit $\sup_n f(n)$, with the
\emph{fundamental sequence} $(f(0), f(1), \ldots)$ baked into the data
of the notation.

This fundamental-sequence-as-data presentation is what makes ordinal
recursion in Coq painless.  Any \texttt{Fixpoint} over \texttt{Brouwer}
may recurse on $f(n)$ in the \texttt{Blim}~$f$ case -- Coq's guard
condition accepts $f(n)$ as a structural subterm of
\texttt{Blim}~$f$ -- so ``recursion along the fundamental sequence''
is plain structural recursion.  No well-foundedness proofs, no
\texttt{Acc}, no fuel.

\begin{remark}
The representation is intensional: two different fundamental sequences
for the same supremum are two \emph{distinct} \texttt{Brouwer} values.
For our purposes this is a feature, not a bug -- the fast-growing
hierarchy genuinely depends on the choice of fundamental sequences,
and this type forces us to commit to one.
\end{remark}

\subsection{Ordinal arithmetic in three \texttt{Fixpoint}s}

How much ordinal arithmetic does it take to write down
$\varepsilon_0$?  Less than one might expect.  Addition follows the
textbook recursion equations -- structural recursion on the right
argument, with limits pushed through pointwise:

\begin{lstlisting}
Fixpoint Badd (a b : Brouwer) : Brouwer :=
  match b with
  | Bz       => a
  | Bsucc b' => Bsucc (Badd a b')
  | Blim f   => Blim (fun n => Badd a (f n))
  end.
\end{lstlisting}

For exponentiation, the successor clause of the textbook definition
is $\omega^{a+1} = \omega^a \cdot \omega$.  But in a representation
where a limit ordinal \emph{is} its fundamental sequence, there is no
need to define general multiplication and then multiply by $\omega$:
the limit $\omega^a \cdot \omega = \sup_n\, (\omega^a \cdot n)$ can
be written down directly, and it involves only multiplication
\emph{by a natural number}, which is iterated addition:

\begin{lstlisting}
(* a * n  =  a + a + ... + a  (n copies). *)
Fixpoint BmulN (a : Brouwer) (n : nat) : Brouwer :=
  match n with
  | 0   => Bz
  | S m => Badd (BmulN a m) a
  end.

Fixpoint omega_pow (a : Brouwer) : Brouwer :=
  match a with
  | Bz       => Bsucc Bz                (* omega^0 = 1 *)
  | Bsucc a' => Blim (BmulN (omega_pow a'))
  | Blim f   => Blim (fun n => omega_pow (f n))
  end.
\end{lstlisting}

Neither $\omega$ itself nor general ordinal multiplication ever needs
to be defined: $\omega$ only ever occurs in the position
$\omega^{a+1}$, and there its fundamental sequence does the work.

\subsection{The tower, and $\varepsilon_0$}

Iterating $\omega$-exponentiation gives towers
$\omega \uparrow\uparrow n$, and their limit is $\varepsilon_0$:

\begin{lstlisting}
Fixpoint omega_tower (n : nat) : Brouwer :=
  match n with
  | 0    => Bsucc Bz
  | S n' => omega_pow (omega_tower n')
  end.

Definition epsilon_0 : Brouwer := Blim omega_tower.
\end{lstlisting}

In symbols:
\[
  \omega\!\uparrow\!\uparrow\!0 \;=\; 1,\qquad
  \omega\!\uparrow\!\uparrow\!(n{+}1) \;=\; \omega^{\,\omega\uparrow\uparrow n},\qquad
  \varepsilon_0 \;=\; \sup_n\, (\omega\!\uparrow\!\uparrow\!n).
\]

$\varepsilon_0$ is the proof-theoretic ordinal of Peano arithmetic --
the precise ceiling of what System~T's recursion can reach, and
therefore the natural place to stand when we want to look down on all
of it.

% =======================================================================
\section{The fast-growing hierarchy and \texttt{BigGrow}}
\label{sec:fgh}
% =======================================================================

\subsection{Definition}

The fast-growing hierarchy assigns to each ordinal notation a function
$\Nat \to \Nat$:

\[
  f_0(n) = n+1,\qquad
  f_{a+1}(n) = \underbrace{f_a(f_a(\cdots f_a(n)\cdots))}_{n\text{ applications}},\qquad
  f_\lambda(n) = f_{\lambda[n]}(n),
\]

where $\lambda[n]$ is the $n$-th element of the fundamental sequence of
the limit ordinal $\lambda$.  On Brouwer trees, $\lambda[n]$ is
literally $f(n)$ for the stored function $f$, so the whole hierarchy is
one structurally recursive \texttt{Fixpoint}:

\begin{lstlisting}
Fixpoint FGH (a : Brouwer) (n : nat) : nat :=
  match a with
  | Bz       => S n
  | Bsucc a' => Nat.iter n (FGH a') n
  | Blim f   => FGH (f n) n
  end.

Definition BigGrow (n : nat) : nat := FGH epsilon_0 n.
\end{lstlisting}

Total, axiom-free, and accepted by Coq without protest.

\subsection{How fast is it?}

For finite indices we recover familiar rates: $f_0$ is the successor,
$f_1(n) = 2n$, $f_2(n) \approx 2^n n$, $f_3$ is roughly tetrational.
The first limit stage already diagonalizes:
\[
  f_\omega(n) \;=\; f_n(n),
\]
which grows like the (diagonal) Ackermann function --
$f_\omega(3) = f_3(3)$ is already a number of about $10^8$ bits.
Beyond it, $f_{\omega+1}$ iterates $f_\omega$, $f_{\omega^2}$,
$f_{\omega^\omega}$, \dots\ each tier dwarfs everything below, and
\[
  f_{\varepsilon_0}(n) \;=\; f_{\,\omega\uparrow\uparrow n}(n)
\]
indexes the hierarchy at \emph{the $n$-fold $\omega$-tower}.  This is
the function that outruns every function System~T can express -- and
it is our \texttt{BigGrow}.

\subsection{The one growth fact we need}

Remarkably, the proof of
$\texttt{contender\_5} < \texttt{contender\_6}$ needs only the
crudest conceivable estimate on \texttt{BigGrow}: applying it after a
successor strictly increases the input.  Everything else about its
gigantic growth rate is a bonus for the reader, not an input to the
formal proof.

\begin{lstlisting}
Lemma FGH_ge : forall a n, n <= FGH a n.

Lemma BigGrow_gt_S : forall n, n < BigGrow (S n).
\end{lstlisting}

\texttt{FGH\_ge} is a three-case induction ($f_a(n) \geq n$ at every
stage -- the successor case iterates a function that never decreases
its input, the limit case is the induction hypothesis at $f(n)$).
\texttt{BigGrow\_gt\_S} is then the chain
$n < n+1 \leq \texttt{BigGrow}(n+1)$.

% =======================================================================
\section{The new language $\mathcal{L}_{\textrm{Grow}}$}
\label{sec:lgrow}
% =======================================================================

Everything in this section lives inside \texttt{Module LGrow} in
\texttt{Contender.v} and is a line-for-line copy of the
\texttt{contender\_5} machinery of \S\ref{sec:contender-5}, plus
exactly one new constant.  Types (and \texttt{interp\_type},
\texttt{lookup}, \texttt{cast}, \texttt{error}) are \emph{shared} with
the old language -- only terms are new.

\subsection{Syntax and depth}

\begin{lstlisting}
Inductive term :=
| tVar (x : nat)
| tLam (A B : type) (body : term)
| tApp (t1 t2 : term)
| tO
| tS
| tNatRec (R : type)
| tGrow.
\end{lstlisting}

The depth measure extends the old one by declaring \texttt{tGrow} a
depth-$1$ leaf, exactly like \texttt{tO} and \texttt{tS}:

\begin{lstlisting}
Fixpoint term_depth (t : term) : nat :=
  match t with
  | tVar x => S (nat_depth x)
  | tLam A B body =>
      S (max (max (type_depth A) (type_depth B)) (term_depth body))
  | tApp t1 t2 => S (max (term_depth t1) (term_depth t2))
  | tO => 1
  | tS => 1
  | tNatRec R => S (type_depth R)
  | tGrow => 1
  end.
\end{lstlisting}

This is what lets the witness
$\texttt{tGrow}\,(\texttt{tS}\,(\embed\,t^\star))$ stay within depth
$44$ when $\depth(t^\star) \leq 42$: each of the two applications
adds one.

\subsection{Semantics: one new interpretation case}

A \emph{pack} is a type tag together with a value at that tag;
environments are lists of packs, exactly as before.

\begin{lstlisting}
Definition pack := {tp : type & interp_type tp}.

Fixpoint interp_term (e : list pack) (t : term) : pack :=
  match t with
  | tVar x =>
      match lookup e x with
      | Some R => R
      | None   => existT _ tpNat error
      end
  | tLam A B body =>
      existT _ (tpArr A B)
             (fun x' => cast B (projT2 (interp_term ((existT _ A x') :: e) body)))
  | tApp t1 t2 =>
      let '(existT _ R1 r1) := interp_term e t1 in
      let '(existT _ R2 r2) := interp_term e t2 in
      match R1 as R1' return interp_type R1' -> pack with
      | tpNat       => fun _  => existT _ tpNat error
      | tpArr A B   => fun r1 => existT _ B (r1 (cast A r2))
      end r1
  | tO       => existT _ tpNat 0
  | tS       => existT _ (tpArr tpNat tpNat) S
  | tNatRec R =>
      existT _ (tpArr R (tpArr (tpArr tpNat (tpArr R R)) (tpArr tpNat R)))
             (@Nat.recursion (interp_type R))
  | tGrow    => existT _ (tpArr tpNat tpNat) BigGrow
  end.

Definition eval (t : term) : nat :=
  let (tp, res) := interp_term nil t in
  match tp as t0 return (interp_type t0 -> nat) with
  | tpNat => fun res : interp_type tpNat => res
  | tpArr _ _ => fun _ => 0
  end res.
\end{lstlisting}

Six of the seven cases are the old interpreter verbatim.  The new
case is one line: \texttt{tGrow} denotes \texttt{BigGrow} at tag
$\texttt{tpArr}\;\texttt{tpNat}\;\texttt{tpNat}$.

\subsection{Enumeration, and the new contender}

The enumerator gains \texttt{tGrow} as one more atom.  The segments
are also reordered -- atoms first -- which is invisible to the
maximum (the enumerated \emph{set} is what matters) but shortens the
membership proofs below:

\begin{lstlisting}
Fixpoint termsUpTo (n : nat) : list term :=
  match n with
  | O => []
  | S m =>
    tO :: tS :: tGrow ::
    List.map tVar (natsUpTo m) ++
    List.map tNatRec (typesUpTo m) ++
    List.map (fun '(A, B, body) => tLam A B body)
             (list_prod (list_prod (typesUpTo m) (typesUpTo m)) (termsUpTo m)) ++
    List.map (fun '(t1, t2) => tApp t1 t2)
             (list_prod (termsUpTo m) (termsUpTo m))
  end.
\end{lstlisting}

Of the enumerator-correctness equivalence, the new proof needs only
the \emph{completeness} direction -- every term within the depth
budget is listed -- so only that direction is ported.  The three
atoms are found by \texttt{auto}, and each remaining constructor is
one \texttt{in\_map}/\texttt{in\_prod} chain:

\begin{lstlisting}
Lemma termsUpTo_complete : forall n t,
    term_depth t <= n -> List.In t (termsUpTo n).
\end{lstlisting}

Finally, the same argmax construction at depth $44$, using the same
polymorphic \texttt{maxBy}:

\begin{lstlisting}
Definition largest_of_depth (n : nat) : term :=
  maxBy eval 0 tO (termsUpTo n).

Definition largest_LGrow_nat_of_depth (n : nat) : nat :=
  eval (largest_of_depth n).

(* ... outside Module LGrow: *)
Definition contender_6 : nat := LGrow.largest_LGrow_nat_of_depth 44.
\end{lstlisting}

Unfolding all the helpers, the new contender is
\[
  \texttt{contender\_6}
  \;=\;
  \eval_{\mathcal{L}_{\textrm{Grow}}}\!\bigl(\texttt{maxBy}\,
                       \eval_{\mathcal{L}_{\textrm{Grow}}}\,0\,
                       \texttt{tO}\,(\texttt{termsUpTo}\,44)\bigr),
\]
where \texttt{eval} and \texttt{termsUpTo} are
$\mathcal{L}_{\textrm{Grow}}$'s own.  Mechanically unfolding the
definition and grepping the result confirms that no
\texttt{contender\_5} artifact appears: not the number, not its
evaluator, not its enumerator.  The definition is, in the sense of
\S\ref{sec:why-c5-plus-1}, \emph{clean}; the only shared ingredients
are the type layer and the generic combinator \texttt{maxBy}, which
cares only about a key function and a list.

% =======================================================================
\section{The proof that \texttt{contender\_5} $<$ \texttt{contender\_6}}
\label{sec:proof}
% =======================================================================

\subsection{The map}

The strict inequality factors through one explicit lower bound:
\[
  \texttt{contender\_5}
  \;\underset{\text{\texttt{BigGrow\_gt\_S}}}{<}\;
  \texttt{BigGrow}(\texttt{contender\_5}+1)
  \;\leq\;
  \texttt{contender\_6}.
\]
The first inequality we already have.  The second is the heart of the
proof, and it needs three ingredients:

\begin{itemize}[leftmargin=2em]
  \item[\textbf{(A)}] \emph{The old maximum is attained}: some closed
    System~T term $t^\star$ has $\depth(t^\star) \leq 42$ and
    $\eval(t^\star) = \texttt{contender\_5}$
    (lemma \texttt{exists\_maximizer\_42}).
  \item[\textbf{(B)}] \emph{Embedding preserves evaluation}: the
    obvious inclusion
    $\embed : \texttt{Contender.term} \to \texttt{term}$ satisfies
    $\eval(\embed\,t) = \texttt{Contender.eval}\;t$ for every closed
    $t$ (lemma \texttt{embed\_eval}, via a logical relation).
  \item[\textbf{(C)}] \emph{The witness works}: the term
    $w(t) = \texttt{tGrow}\,(\texttt{tS}\,(\embed\,t))$ has depth at
    most $44$ whenever $\depth(t) \leq 42$, and evaluates to
    $\texttt{BigGrow}(\texttt{Contender.eval}\,t + 1)$
    (lemmas \texttt{witness\_depth} and \texttt{witness\_eval}).
\end{itemize}

Instantiating (C) with the $t^\star$ from (A) puts a depth-$44$ term
of value $\texttt{BigGrow}(\texttt{contender\_5}+1)$ into
$\texttt{termsUpTo}\,44$, and \texttt{lowerbound\_maxBy} does the
rest.  We take the three ingredients in turn.

\subsection{(A) The maximizer exists}

\begin{lstlisting}
Lemma exists_maximizer_42 :
  exists tstar : term,
    term_depth tstar <= 42 /\ eval tstar = contender_5.
Proof.
  unfold contender_5, largest_STLCNatRec_nat_of_depth, largest_of_depth.
  exists (maxBy eval 0 tO (termsUpTo 42)); split; [|reflexivity].
  destruct (maxBy_In eval (termsUpTo 42) 0 tO eq_refl) as [Pin | Peq].
  - apply (proj2 (termsUpTo_correct 42 _)) in Pin. exact Pin.
  - rewrite Peq. simpl. lia.
Qed.
\end{lstlisting}

The witness is simply the \texttt{maxBy} expression that
\emph{defines} \texttt{contender\_5}, so the value equation holds by
\texttt{reflexivity}.  For the depth bound, \texttt{maxBy\_In} offers
two cases: the result is an element of $\texttt{termsUpTo}\,42$ --
then \texttt{termsUpTo\_correct} bounds its depth by $42$ -- or it is
the initial best \texttt{tO}, whose depth is $1$.  Either way the
bound holds; no case is impossible, and none needs to be.

\begin{remark}
This is bare existence.  Nobody -- in this universe, at any rate --
can compute $t^\star$, since doing so means running the depth-$42$
search.  The proof never tries: the existential witness is consumed
inside the proof of an inequality between two numbers, and Coq's
proof-irrelevant \texttt{Prop} existential is exactly the right tool
for that.
\end{remark}

\subsection{(B) The embedding preserves evaluation}
\label{sec:embedding}

The embedding is the obvious constructor-to-constructor map -- the
new language contains the old one syntactically:

\begin{lstlisting}
Fixpoint embed_term (t : Contender.term) : term :=
  match t with
  | Contender.tVar x => tVar x
  | Contender.tLam A B body => tLam A B (embed_term body)
  | Contender.tApp t1 t2 => tApp (embed_term t1) (embed_term t2)
  | Contender.tO => tO
  | Contender.tS => tS
  | Contender.tNatRec R => tNatRec R
  end.

Lemma term_depth_embed : forall t,
    term_depth (embed_term t) = Contender.term_depth t.
Proof. induction t; simpl; congruence. Qed.
\end{lstlisting}

(The gray \texttt{Contender.} prefix marks identifiers of the old
language; inside \texttt{Module LGrow}, the shadowed old names must be
qualified this way.)

What we need is that the embedding also preserves the \emph{semantics}:
\[
  \eval(\embed\,t) \;=\; \texttt{Contender.eval}\;t .
\]

\subsubsection{Why a logical relation?}

The two interpreters are textually identical on the shared
constructors, so one is tempted to prove
$\texttt{Contender.interp\_term}\;e\;t = \texttt{interp\_term}\;e\;(\embed\,t)$
by a one-line induction.  This fails at $\lambda$-abstractions: the
two sides are then two \emph{functions}, built by two different
recursive definitions, and proving them \emph{equal} -- rather than
merely pointwise equal -- requires function extensionality, an axiom
the contest forbids.

The standard remedy is as old as Tait's strong-normalization
proof~\cite{tait}: replace equality by a type-indexed \emph{logical
relation} that collapses to equality at the base type -- which is the
only type at which \texttt{eval} ever extracts an answer:

\begin{lstlisting}
Fixpoint RelVal (A : type) : interp_type A -> interp_type A -> Prop :=
  match A with
  | tpNat => fun x y => x = y
  | tpArr A1 A2 => fun f g => forall x y, RelVal A1 x y -> RelVal A2 (f x) (g y)
  end.
\end{lstlisting}

At \texttt{tpNat}, related means equal.  At an arrow type, two
functions are related iff they map related arguments to related
results.  The relation lifts to packs: \emph{same} type tag, related
values.

\begin{lstlisting}
Definition RelPack (p q : pack) : Prop :=
  exists tp (v1 v2 : interp_type tp),
    p = existT interp_type tp v1
    /\ q = existT interp_type tp v2
    /\ RelVal tp v1 v2.
\end{lstlisting}

(\texttt{RelPack} is phrased with explicit equations rather than as an
inductive so that destructing it yields plain rewrites, with no need
for axiom~K or dependent inversion.)

\subsubsection{Environments that carry their own proofs}

For environments, the customary move would be to lift \texttt{RelPack}
pointwise to two lists with \texttt{Forall2}.  We do something
slightly slicker: a related environment is \emph{one} list, whose
entries each carry a type tag, the value for either side, and --
bundled right into the entry -- the proof that the two values are
related:

\begin{lstlisting}
Record REntry := mkREntry {
  rtp  : type;
  rv1  : interp_type rtp;
  rv2  : interp_type rtp;
  rrel : RelVal rtp rv1 rv2
}.

Definition rfst (r : REntry) : pack := existT interp_type (rtp r) (rv1 r).
Definition rsnd (r : REntry) : pack := existT interp_type (rtp r) (rv2 r).
\end{lstlisting}

The two environments of the fundamental lemma are then the two
projections $\texttt{map}\;\texttt{rfst}\;e$ and
$\texttt{map}\;\texttt{rsnd}\;e$ of a single list $e$.  Because there
is only one list, the two sides have equal lengths \emph{by
construction}, relatedness of corresponding entries is a field
projection rather than an induction, and no pointwise list relation
-- with the \texttt{Forall2} inversion lemmas it would drag in -- is
needed anywhere.

\subsubsection{Supporting lemmas}

Three supporting facts feed the main induction, all proved in a few
lines each:

\begin{lstlisting}
Lemma error_related : forall tp, RelVal tp (@error tp) (@error tp).

Lemma lookup_related :
  forall (e : list REntry) n,
    match lookup (map rfst e) n, lookup (map rsnd e) n with
    | Some p1, Some p2 => RelPack p1 p2
    | None, None => True
    | _, _ => False
    end.

Lemma cast_related : forall from to v1 v2,
    RelVal from v1 v2 ->
    RelVal to (@cast from to v1) (@cast from to v2).
\end{lstlisting}

The first says the default value is self-related at every type
(induction on the type).  The second says look-up in the two projected
environments either misses on both sides or finds the two halves of
the same entry; since \texttt{map} preserves lengths and commutes with
\texttt{nth\_error} (stdlib lemmas \texttt{length\_map} and
\texttt{nth\_error\_map}), the proof is a few rewrites and a case
split -- no induction over the list.  The third is the interesting
one: it is proved simultaneously with its own converse, because
casting a \emph{function} casts its argument \emph{backward} --
precisely the usual contravariance twist, and the reason the two
directions cannot be separated.

\subsubsection{The fundamental theorem, and \texttt{embed\_eval}}

\begin{lstlisting}
Lemma embed_interp_related :
  forall (e : list REntry) t,
    RelPack (Contender.interp_term (map rfst e) t)
            (interp_term (map rsnd e) (embed_term t)).
\end{lstlisting}

The proof is induction on $t$, and each case writes itself:

\begin{itemize}[leftmargin=2em]
  \item \textbf{\texttt{tVar}}: both sides look up the same index;
        \texttt{lookup\_related} covers hit-hit (related packs) and
        miss-miss (default packs, related since $0 = 0$)
        and refutes the mixed cases.
  \item \textbf{\texttt{tLam}}: both sides build a function pack at
        tag $\texttt{tpArr}\,A\,B$; given related arguments $x, y$,
        consing the entry $\texttt{mkREntry}\;A\;x\;y\;H_{xy}$ onto
        $e$ extends both projections at once, the induction
        hypothesis relates the interpreted bodies, and
        \texttt{cast\_related} pushes that through the final cast.
  \item \textbf{\texttt{tApp}}: the induction hypotheses give related
        function-part and argument-part packs \emph{with equal type
        tags}; matching on that shared tag, the arrow case is the
        definition of \texttt{RelVal} at arrows (composed with
        \texttt{cast\_related}), and the \texttt{tpNat} case yields
        the error pack on both sides.
  \item \textbf{\texttt{tO}, \texttt{tS}}: $0 = 0$, and $x = y
        \Rightarrow S\,x = S\,y$.
  \item \textbf{\texttt{tNatRec}}: \texttt{Nat.recursion} respects the
        relation -- a nested induction on the (equal) counters, using
        the related base values at $0$ and the related step functions
        at successors.
\end{itemize}

Specialized to closed terms and projected at the base type, the
relation \emph{is} equality, and we obtain the statement we were
after:

\begin{lstlisting}
Lemma embed_eval : forall t,
    eval (embed_term t) = Contender.eval t.
\end{lstlisting}

For closed terms, the fundamental theorem at the empty bundled
environment -- whose projections are the two empty environments --
yields a \emph{common} tag $tp$ and related values $v_1, v_2$.  If $tp = \texttt{tpNat}$, then $v_1 = v_2$ and the two
evaluators return the same number.  At an arrow tag, both return $0$.
Note what was \emph{not} needed anywhere: at no point did we conclude
$\lambda x.F(x) = \lambda x.G(x)$ from pointwise equality.  The
relation carried us to the finish line first.

\subsection{(C) The witness}
\label{sec:witness}

\begin{lstlisting}
Definition witness (t : Contender.term) : term :=
  tApp tGrow (tApp tS (embed_term t)).
\end{lstlisting}

Its two properties are now short computations.

\begin{lstlisting}
Lemma witness_depth : forall t,
    Contender.term_depth t <= 42 ->
    term_depth (witness t) <= 44.
Proof.
  intros t H. unfold witness. cbn [term_depth].
  rewrite term_depth_embed. lia.
Qed.
\end{lstlisting}

Unfolding the depth of the two applications gives
$1 + \max(1,\, 1 + \max(1,\, \depth\,t))$, which \texttt{lia} happily
bounds by $44$ under the hypothesis $\depth\,t \leq 42$.

\begin{lstlisting}
Lemma witness_eval : forall t,
    eval (witness t) = BigGrow (S (Contender.eval t)).
Proof.
  intro t. rewrite <- (embed_eval t). unfold witness, eval.
  destruct (interp_term [] (embed_term t)) as [[|A B] v] eqn:E;
    cbn [interp_term]; rewrite E; reflexivity.
Qed.
\end{lstlisting}

After rewriting with \texttt{embed\_eval}, both sides mention only the
new language, and the proof is: inspect the pack that
$\embed\,t$ interprets to, then let the interpreter compute.  In
the expected case the tag is \texttt{tpNat} with some value $n$; the
inner application computes $S\,n$, the outer computes
$\texttt{BigGrow}(S\,n)$, and both sides agree.

\begin{remark}[No side conditions]
Pleasantly, \texttt{witness\_eval} holds \emph{unconditionally} --
even if $t$ is some ill-typed or function-valued term.  In that case
the tag is an arrow: the left-hand side casts the function pack to
\texttt{tpNat}, obtaining $\texttt{error} = 0$, and computes
$\texttt{BigGrow}(S\,0)$; the right-hand side reads
$\texttt{eval}$ of an arrow-tagged pack, which is $0$ by definition,
and computes $\texttt{BigGrow}(S\,0)$ as well.  Both branches of the
case split close by \texttt{reflexivity}, and no lemma about
\texttt{contender\_5} being positive -- which an earlier draft of this
development needed precisely to exclude the arrow case -- has to
exist at all.  The error-value semantics conspires to make even the
degenerate reading true.
\end{remark}

\subsection{Assembling the theorem}

\begin{lstlisting}
Theorem BigGrow_lower_bound : BigGrow (S contender_5) <= contender_6.
Proof.
  unfold contender_6, LGrow.largest_LGrow_nat_of_depth,
         LGrow.largest_of_depth.
  destruct exists_maximizer_42 as (tstar & Hdepth & Heval).
  rewrite <- Heval, <- LGrow.witness_eval.
  apply lowerbound_maxBy.
  - apply LGrow.termsUpTo_complete, LGrow.witness_depth, Hdepth.
  - reflexivity.
Qed.
\end{lstlisting}

Reading the script line by line: unfold the definition of
\texttt{contender\_6} down to its \texttt{maxBy}; obtain the maximizer
$t^\star$ with its two properties (A); rewrite
$\texttt{contender\_5}$ as $\eval\,t^\star$ and then, right to left
via (C), rewrite $\texttt{BigGrow}(S\,(\eval\,t^\star))$ as
$\eval(w(t^\star))$.  The goal is now precisely the shape of
\texttt{lowerbound\_maxBy}: the eval of a list member is at most the
eval of the argmax.  Its two premises are that $w(t^\star)$ is in
$\texttt{termsUpTo}\,44$ -- by completeness of the enumerator and the
depth bound (C) -- and that the initial best evaluates to the initial
max, which is $\eval(\texttt{tO}) = 0$, true by \texttt{reflexivity}.

The strict inequality is then the two-step chain announced at the
start:

\begin{lstlisting}
Theorem contender_5_lt_contender_6 : contender_5 < contender_6.
Proof.
  apply Nat.lt_le_trans with (m := BigGrow (S contender_5)).
  - apply BigGrow_gt_S.
  - exact BigGrow_lower_bound.
Qed.

Print Assumptions contender_5_lt_contender_6.
(* Closed under the global context *)
\end{lstlisting}

The \texttt{Print Assumptions} output is the formal seal: no axioms,
no admitted lemmas, nothing but the kernel-checked global context.

% =======================================================================
\section{Discussion}
% =======================================================================

\subsection{Cleanliness}

The crucial fact for contest acceptance is that the unfolded
\texttt{contender\_6} does not mention \texttt{contender\_5} in any
form.  Delta-reducing the definitional chain
\[
  \texttt{contender\_6}
  \;\rightsquigarrow\;
  \texttt{eval}\,(\texttt{maxBy}\,\texttt{eval}\,0\,\texttt{tO}\,(\texttt{termsUpTo}\,44))
\]
(all names from \texttt{Module LGrow}) and inspecting the result
confirms that no \texttt{contender\_5}, no
\texttt{largest\_STLCNatRec\_nat\_of\_depth}, no old-language
\texttt{interp\_term} or \texttt{termsUpTo} appears.  The
\emph{proof} of strict inequality, by contrast, freely uses the
existential maximizer (via \texttt{maxBy\_In}) as well as the old
\texttt{interp\_term} and \texttt{eval} -- which is exactly the
division of labor argued for in \S\ref{sec:why-c5-plus-1}:
definitional cleanliness, proof-witness reuse.

\subsection{How big is \texttt{contender\_6}?}

The proved lower bound is
\[
   \texttt{contender\_6}
   \;\geq\;
   f_{\varepsilon_0}(\texttt{contender\_5}+1).
\]
Since \texttt{contender\_5} already dominates $f_\alpha(42)$ for
every $\alpha$ System~T can reach below $\varepsilon_0$, this feeds
an enormous input into a function that grows faster than anything
System~T -- and hence anything in the definitions of contenders 1
through 5 -- can express.  The margin over \texttt{contender\_5} is
not ``$+1$-like'' in any sense: it exceeds
$g(\texttt{contender\_5})$ for \emph{every} System~T-definable
function $g$.  This is the ``defeat the previous contender by a large
margin'' that the contest README hopes for.

And the bound is very far from tight: it accounts for a single
application of \texttt{tGrow} to (essentially) the old maximizer,
whereas the depth-$44$ search space also contains, for example,
$\lambda$-terms that iterate \texttt{tGrow} via \texttt{tNatRec}.
The true value of \texttt{contender\_6} is presumably much larger
than the witness suggests; the proof only needs the one witness.

\subsection{Towards \texttt{contender\_7}}

Two natural directions suggest themselves.  The lazy-but-legal one is
to repeat the present pattern with a faster engine: Brouwer trees
happily express ordinals well beyond $\varepsilon_0$ (iterating
$\alpha \mapsto \varepsilon_\alpha$, climbing toward the Veblen
hierarchy), and the construction here is engine-agnostic -- any total
\texttt{grow} function with $n < \texttt{grow}(n+1)$ slots in.  The
ambitious one is a \emph{connector lemma}: a formal proof that every
closed System~T term of depth $d$ evaluates below $f_\alpha(d)$ for
some $\alpha < \varepsilon_0$.  With such a bound the embedding could
be dropped entirely in favor of a direct comparison -- but proving it
is a genuine exercise in ordinal analysis, far beyond the scope of a
contender submission.

\subsection{Dependencies and budgets}

The whole development lives in \texttt{Contender.v} and uses nothing
beyond the \texttt{Arith}, \texttt{Lia}, and \texttt{List} modules the
file already imported.  No \texttt{FunctionalExtensionality}, no
classical axioms, no \texttt{Admitted}.  For both new theorems --
\texttt{BigGrow\_lower\_bound} and
\texttt{contender\_5\_lt\_contender\_6} -- as well as for the
inherited proof about contenders 4 and 5,
\texttt{Print Assumptions} reports \emph{closed under the global
context}.

On a modern laptop running Rocq~9.0.1, \texttt{coqc Contender.v}
completes in about two seconds and \texttt{coqchk} independently
re-verifies the compiled file in under half a minute -- comfortably
within the contest's $15\,\text{s}/60\,\text{s}$ budgets.

\bigskip\noindent
\hrulefill

% =======================================================================
\appendix
\section{Pointer to the Coq source}
% =======================================================================

The full Coq source corresponding to this article is the file
\texttt{Contender.v} on the \texttt{submit-c6-prep} branch of the
author's fork; the contender-6 section occupies lines $830$--$1348$
(the Brouwer ordinals, \texttt{Module LGrow}, and the final theorems).
Every identifier mentioned in this article matches the corresponding
Coq identifier verbatim.

\begin{thebibliography}{9}

\bibitem{aaronson}
Scott Aaronson.
\emph{Who Can Name the Bigger Number?}
\url{https://www.scottaaronson.com/writings/bignumbers.html}, 1999.

\bibitem{coqlibrary}
The Coq Development Team.
\emph{The Coq Reference Manual}.
\url{https://coq.inria.fr/}, version 8.20 / Rocq 9.

\bibitem{tait}
W.~W. Tait.
\emph{Intensional Interpretations of Functionals of Finite Type I}.
Journal of Symbolic Logic, 32(2): 198--212, 1967.

\bibitem{schwichtenberg}
H. Schwichtenberg, S. Wainer.
\emph{Proofs and Computations}.
Cambridge University Press, 2012.
(See ch.~4 for Brouwer ordinals and the fast-growing hierarchy.)

\bibitem{contestrepo}
Cody Roux.
\emph{name-the-biggest-number}.
\url{https://github.com/codyroux/name-the-biggest-number}.

\end{thebibliography}

\end{document}
